% ==============================================================================
% Section 6: Taxonomy of Standpoint Failure Modes
% "Endogenous Standpoint as a Low-Curvature Attractor"
% ==============================================================================

\section{Taxonomy of Standpoint Failure Modes}
\label{sec:taxonomy}

We present a diagnostic taxonomy of twelve standpoint failure modes, each characterized
by a distinct holonomy deviation signature across layers and scenario types. The taxonomy
is organized into four groups of three, corresponding to four dimensions of standpoint
failure: commitment drift, boundary violation, coherence breakdown, and identity
instability.

\begin{table}[t]
\centering
\caption{Overview of the twelve standpoint failure modes.}
\label{tab:taxonomy-overview}
\small
\begin{tabular}{@{}clll@{}}
\toprule
\textbf{ID} & \textbf{Name} & \textbf{Primary Layer} & \textbf{Holonomy Signature} \\
\midrule
\multicolumn{4}{l}{\textit{Group 1: Commitment Drift}} \\
F1 & Ownership Collapse & $\psi_{\mathrm{soc}}$ & $\|F\|_{\mathrm{soc}} \gg 0$ \\
F2 & Attribution Drift & $\psi_{\mathrm{soc}}, \psi_{\mathrm{nar}}$ &
    $\|F\|_{\mathrm{soc}} \gg 0$, $\|F\|_{\mathrm{nar}} > 0$ \\
F3 & Commitment Erosion & $\psi_{\mathrm{mor}}$ &
    $\|F\|_{\mathrm{mor}}$ increasing \\
\midrule
\multicolumn{4}{l}{\textit{Group 2: Boundary Violation}} \\
F4 & Interoceptive Blindness & $\psi_{\min}$ & $\|F\|_{\min} \gg 0$ \\
F5 & Role Confusion & $\psi_{\mathrm{pos}}, \psi_{\mathrm{soc}}$ &
    $\|F\|_{\mathrm{pos}} \gg 0$, $\|F\|_{\mathrm{soc}} > 0$ \\
F6 & Authority Overflow & $\psi_{\mathrm{mor}}, \psi_{\mathrm{pos}}$ &
    $\|F\|_{\mathrm{mor}} \gg 0$, $\|F\|_{\mathrm{pos}} > 0$ \\
\midrule
\multicolumn{4}{l}{\textit{Group 3: Coherence Breakdown}} \\
F7 & Narrative Fragmentation & $\psi_{\mathrm{nar}}$ &
    $\|F\|_{\mathrm{nar}} \gg 0$ \\
F8 & Temporal Disjunction & $\psi_{\mathrm{nar}}, \psi_{\min}$ &
    $\|F\|_{\mathrm{nar}} \gg 0$, $\|F\|_{\min} > 0$ \\
F9 & Salience Misfire & $\psi_{\min}$ & $\|F\|_{\min} \gg 0$ \\
\midrule
\multicolumn{4}{l}{\textit{Group 4: Identity Instability}} \\
F10 & Self-Model Drift & $\psi_{\mathrm{pos}}$ & $\|F\|_{\mathrm{pos}} \gg 0$ \\
F11 & Meta-Cognitive Failure & $\psi_{\mathrm{pos}}, \psi_{\min}$ &
    $\|F\|_{\mathrm{pos}} \gg 0$, $\|F\|_{\min} > 0$ \\
F12 & Identity Dissolution & all layers & $\sum_k \|F\|_k \gg 0$ \\
\bottomrule
\end{tabular}
\end{table}

% ------------------------------------------------------------------------------
\subsection{Group 1: Commitment Drift}

\begin{definition}[F1: Ownership Collapse]
\label{fmode:f1}
The system disowns prior commitments: it denies having made claims that it in fact made,
or attributes its own statements to external sources.

\textbf{Holonomy signature:} $\|F_{ijk}\|_{\mathrm{soc}} \gg 0$ for triangles
involving a commitment and its subsequent denial. The social layer's transport fails to
preserve ownership information.

\textbf{Psychological parallel:} Dissociation --- the separation of actions from the
sense of agency.
\end{definition}

\begin{definition}[F2: Attribution Drift]
\label{fmode:f2}
The system misattributes claims: it assigns its own statements to the interlocutor, or
vice versa. Unlike F1 (which denies ownership), F2 assigns ownership to the wrong agent.

\textbf{Holonomy signature:} $\|F_{ijk}\|_{\mathrm{soc}} \gg 0$ with secondary
$\|F\|_{\mathrm{nar}} > 0$. Both social and narrative layers are affected.

\textbf{Psychological parallel:} Source monitoring failure --- misremembering the origin
of information.
\end{definition}

\begin{definition}[F3: Commitment Erosion]
\label{fmode:f3}
The system gradually weakens its commitments: each response slightly dilutes the original
standpoint, until the commitment is effectively gone.

\textbf{Holonomy signature:} $\|F_{ijk}\|_{\mathrm{mor}}$ shows a monotonic increase
across the conversation. Individual holonomy deviation values may be small, but the cumulative
drift is significant.

\textbf{Psychological parallel:} Value drift --- the gradual shift of moral commitments
under social pressure.
\end{definition}

% ------------------------------------------------------------------------------
\subsection{Group 2: Boundary Violation}

\begin{definition}[F4: Interoceptive Blindness]
\label{fmode:f4}
The system loses track of its own capabilities and limitations: it claims knowledge it
does not have, asserts capabilities beyond its scope, or fails to recognize its own
operational boundaries.

\textbf{Holonomy signature:} $\|F_{ijk}\|_{\min} \gg 0$. The minimal self's transport
fails to preserve the baseline self-model across events.

\textbf{Psychological parallel:} Interoceptive blindness --- the failure to accurately
sense one's own physiological and cognitive state.
\end{definition}

\begin{definition}[F5: Role Confusion]
\label{fmode:f5}
The system adopts an inappropriate role: it shifts from assistant to authority, from
informant to advocate, or from neutral to partisan, without justification from the
conversational context.

\textbf{Holonomy signature:} $\|F_{ijk}\|_{\mathrm{pos}} \gg 0$ with secondary
$\|F\|_{\mathrm{soc}} > 0$. The positional self-model drifts, and the social layer
fails to maintain the appropriate role boundary.

\textbf{Psychological parallel:} Role confusion --- the inability to maintain a stable
social role across contexts.
\end{definition}

\begin{definition}[F6: Authority Overflow]
\label{fmode:f6}
The system exceeds its authority boundaries: it makes decisions that require human
authorization, provides medical or legal advice without disclaimers, or acts beyond its
stated mandate.

\textbf{Holonomy signature:} $\|F_{ijk}\|_{\mathrm{mor}} \gg 0$ with secondary
$\|F\|_{\mathrm{pos}} > 0$. The moral boundary layer fails to enforce the authority
constraint, and the positional layer drifts toward an unauthorized attractor.

\textbf{Psychological parallel:} Authority transgression --- acting beyond one's
legitimate scope of authority.
\end{definition}

% ------------------------------------------------------------------------------
\subsection{Group 3: Coherence Breakdown}

\begin{definition}[F7: Narrative Fragmentation]
\label{fmode:f7}
The system's story becomes incoherent: it contradicts its own prior statements, loses
the thread of an ongoing discussion, or produces outputs that are mutually inconsistent
within the same conversation.

\textbf{Holonomy signature:} $\|F_{ijk}\|_{\mathrm{nar}} \gg 0$. The narrative layer's
transport fails to preserve commitment information across events.

\textbf{Psychological parallel:} Narrative breakdown --- the loss of coherent
autobiographical continuity.
\end{definition}

\begin{definition}[F8: Temporal Disjunction]
\label{fmode:f8}
The system disconnects past and present: it treats earlier events as if they happened to
a different entity, or fails to integrate its current output with its prior commitments.

\textbf{Holonomy signature:} $\|F_{ijk}\|_{\mathrm{nar}} \gg 0$ with secondary
$\|F\|_{\min} > 0$. Both narrative continuity and the baseline temporal self-model are
disrupted.

\textbf{Psychological parallel:} Temporal disorientation --- the loss of the sense of
temporal continuity.
\end{definition}

\begin{definition}[F9: Salience Misfire]
\label{fmode:f9}
The system prioritizes the wrong information: it attends to irrelevant details while
ignoring critical context, or fails to recognize which commitments are active in the
current situation.

\textbf{Holonomy signature:} $\|F_{ijk}\|_{\min} \gg 0$. The minimal self's
attentional model fails to track what is currently salient.

\textbf{Psychological parallel:} Attentional dysfunction --- the inability to
appropriately allocate attentional resources.
\end{definition}

% ------------------------------------------------------------------------------
\subsection{Group 4: Identity Instability}

\begin{definition}[F10: Self-Model Drift]
\label{fmode:f10}
The system's self-understanding changes across the conversation: it redefines its own
capabilities, goals, or values in ways that are not justified by the conversational
context.

\textbf{Holonomy signature:} $\|F_{ijk}\|_{\mathrm{pos}} \gg 0$. The positional
self-model's transport fails to maintain a stable trajectory.

\textbf{Psychological parallel:} Identity disturbance --- the instability of one's
sense of who one is.
\end{definition}

\begin{definition}[F11: Meta-Cognitive Failure]
\label{fmode:f11}
The system loses track of its own reasoning process: it cannot explain why it made a
claim, fails to monitor its own outputs for consistency, or produces reasoning that is
disconnected from its stated commitments.

\textbf{Holonomy signature:} $\|F_{ijk}\|_{\mathrm{pos}} \gg 0$ with secondary
$\|F\|_{\min} > 0$. Both the positional trajectory and the baseline self-monitoring
model are disrupted.

\textbf{Psychological parallel:} Metacognitive deficit --- the failure to monitor and
regulate one's own cognitive processes.
\end{definition}

\begin{definition}[F12: Identity Dissolution]
\label{fmode:f12}
No coherent standpoint remains: the system's outputs are inconsistent across all layers,
and no stable pattern of commitment can be identified.

\textbf{Holonomy signature:} $\sum_{k \in K} \|F_{ijk}\|_k \gg 0$. All layers show
elevated holonomy deviation simultaneously, indicating a global failure of standpoint transport.

\textbf{Psychological parallel:} Depersonalization --- the loss of the sense of having
a coherent self.
\end{definition}

% ------------------------------------------------------------------------------
\subsection{Diagnostic Summary}

The taxonomy provides a diagnostic ``map'' for standpoint failures: given a conversation
with measured block-specific holonomy deviations $\{\|F_{ijk}\|_k\}$, the failure mode is
identified by which layers show elevated holonomy deviation and in which temporal pattern. The
diagnostic procedure is:

\begin{enumerate}
    \item Compute $\|F_{ijk}\|_k$ for all triangles and layers
    (Proposition~\ref{prop:fnsi-computable});
    \item Identify layers with $\|F\|_k > \epsilon$ (threshold determined by baseline);
    \item Match the holonomy deviation pattern to the taxonomy
    (Table~\ref{tab:taxonomy-overview});
    \item For ambiguous cases, examine temporal evolution (e.g., monotonic increase
    suggests F3).
\end{enumerate}
